\section{Open problems}\label{sec:outlook}

\emph{More points per cell.} The defect of Definition~\ref{def:D} is defined at every point of every locally finite configuration. For
a periodic configuration $X$ with $N$ points $x_1,\dots,x_N$ per cell let $\overline D(X):=\frac1N\sum_iD(X,x_i)$.

\begin{conjecture}\label{conj:all}
There is a constant $c>0$ such that $e(X)-\eB\ge c\,\overline D(X)$ for every periodic configuration $X\subset\R^3$ of density
$3/(4\pi)$.
\end{conjecture}

This would imply that the BCC lattice is the unique periodic minimizer of the Coulomb energy at fixed density. Indeed, if
$\overline D(X)=0$, then the neighbours of every point of $X$ within $\frac{13}{10}\ltwo$ form a rotated copy of the star
$\ltwo\Star$; if $y$ is one of the eight nearest neighbours of $x$, then $x$ and three of the six second neighbours of $x$ are
nearest neighbours of $y$, which forces the star at $y$ to coincide with the star at $x$; since the nearest-neighbour vectors
of the star generate the BCC lattice, $X$ contains a BCC lattice $\Lambda$ of the same density; and every point of $X$ lies within
the covering radius $\frac{\sqrt5}4\ltwo<\frac{13}{10}\ltwo$ of $\Lambda$ from a point of $\Lambda$, hence is a neighbour of that
point and belongs to $\Lambda$. The conjecture is also the form
of the statement that is useful for the analysis of low-energy configurations that are not periodic. Theorem~\ref{thm:main} proves the conjecture for $N\le2$ with $c=\gap$, and the FCC lattice shows that
$c\le\gap$ in general. We do not know whether $c=\gap$ is admissible for all $N$. The methods of this paper extend in principle to
any fixed $N$: the positive series \eqref{eq:seriesN}, the compactness argument and the vertex principle are available, and the
anchor mechanism of Section~\ref{ssec:anchor} applies to a neighbourhood of the configurations with fewer points per cell. The
dimension of the configuration space, $5+3(N-1)$, makes an exhaustive cover expensive already for $N=3$, and a proof for all $N$
needs a different idea.

\emph{Other interactions.} The representation by a series with nonnegative terms holds for the Epstein zeta function at all
exponents below the pole \cite{SarnakStrombergsson2006}, where the optimal lattice is expected to change from BCC to FCC, and a
similar transition has been established among lattices, with computer assistance, for Yukawa interactions \cite{LuoWei2026Yukawa}. It would be interesting to know whether the
two-site analogue of Theorem~\ref{thm:min} holds throughout these families, with the minimizer among lattices also minimizing among
two-site configurations, or whether there are parameters at which a genuinely two-site structure wins.

\emph{Formal verification.} A certificate of the two-site cover with fewer and larger boxes, built on sharper bounds near the
set of lattices and designed for evaluation in the Lean kernel, as was done for the lattice case, together with formal checking
functions for the remaining kinds of rules (Section~\ref{ssec:lean}), would make Theorem~\ref{thm:main} a formally verified
theorem.

\emph{A conceptual proof.} For lattices in two dimensions the optimality of the triangular lattice follows from Montgomery's theorem
on theta functions, by an argument that covers all completely monotone interactions at once. In three dimensions the optimal
lattice depends on the interaction, and the Coulomb energy of the FCC lattice exceeds that of the BCC lattice by six parts in
$10^5$; we do not expect an argument that avoids quantitative information of this precision, but it would be very desirable to
replace the cover by a structural argument, at least in the neighbourhood of the set of lattices where the anchor rules are used.

\section*{Declaration on the use of AI systems}
The mathematical arguments, the programs, the certificates, the formalization and the text of this paper were produced with
extensive use of AI systems, namely large language models of Anthropic (Claude, including Claude Opus~5.5 and Claude Fable~5.2) and
of OpenAI (used through Codex and recorded as Astra and Sol~6.1--6.2), working under the direction of the author. The checks
described in Section~\ref{ssec:checks}, and the corresponding checks of the formal development, were likewise performed by
instances of these systems, in sessions separate from those that produced the material checked. The validity of the results does not rest on these checks: it
rests on the proofs given in this paper, on the certificates and the verifying programs, and, for the parts stated in
Section~\ref{ssec:lean}, on the Lean kernel.
